% GENERATED FILE. Edit canonical lessons, then run npm run generate:books.
% canonical-source-hash: fnv1a64:8cff8ca0547938a7
% canonical-lessons: ZH-C142-tuodi, ZH-C142-laji, ZH-C142-shai, ZH-C142-zhengli, ZH-C142-jiaohua

\chapter{More Housework}
\label{ch:zh-more-housework}
\hlchaptermodality{142}

\begin{chapteropening}
\textbf{By the end of this chapter:} \emph{I can talk about mopping, rubbish, drying things in the sun, sorting things out and watering the plants.}

Complete the last lesson of chapter 142: I can talk about mopping, rubbish, drying things in the sun, sorting things out and watering the plants.
\end{chapteropening}

\section[tuōdì]{tuōdì — to mop the floor}

\label{lesson:ZH-C142-tuodi}



\begin{quote}
Before the new one: say the Mandarin for to take a walk, then the Mandarin for to go fishing.
\end{quote}

\subsection*{You'll want to know: tuōdì}
\emph{tuōdì} — \textquotedblleft{}to mop the floor\textquotedblright{}.

\begin{quote}
\emph{wǒ měitiān tuōdì} — I mop the floor every day.
\end{quote}

Its characters come later; for now it is a word for the ear.

\begin{cousinweb}
\emph{tuōdì} is built from \emph{tuō} 'to drag' and \emph{dì} 'ground', both inherited from Middle Chinese, \emph{tha} and \emph{dijH} (Guangyun, in Baxter's transcription). The dictionary glosses the first 'to drag' under an older way of writing it, and notes that it is popularly written the modern way; \emph{tha} is also the old form of the word for he, a different word.
\end{cousinweb}

\begin{grammarlens}[title={What you've built}]
One more word for everyday actions.
\end{grammarlens}

\subsection*{Your turn}
\begin{itemize}
\raggedright
  \item \emph{Say it:} \emph{tuōdì}
  \item \emph{Say it:} \emph{tuōdì}, once more
  \item \emph{Say it:} \emph{wǒ měitiān tuōdì} — I mop the floor every day.
  \item \emph{Recall:} say the Mandarin for to take a walk, then the Mandarin for to go fishing, then say \emph{tuōdì} again
\end{itemize}

\begin{tcolorbox}[breakable,colback=teal!4,colframe=teal!35!black,title={Before you move on}]
What does tuōdì mean? (\textquotedblleft{}To mop the floor\textquotedblright{}.) Say it once more.
\end{tcolorbox}
\section[lājī]{lājī — rubbish}

\label{lesson:ZH-C142-laji}



\begin{quote}
Before the new one: say the Mandarin for to go fishing, then the Mandarin for to mop the floor.
\end{quote}

\subsection*{You'll want to know: lājī}
\emph{lājī} — \textquotedblleft{}rubbish\textquotedblright{}.

\begin{quote}
\emph{ménkǒu yǒu lājī} — There's rubbish by the door.
\end{quote}

Its characters come later; for now it is a word for the ear.

\begin{cousinweb}
\emph{lājī} No dictionary we use records an older form, so its history is not settled.
\end{cousinweb}

\begin{grammarlens}[title={What you've built}]
One more word for everyday actions.
\end{grammarlens}

\subsection*{Your turn}
\begin{itemize}
\raggedright
  \item \emph{Say it:} \emph{lājī}
  \item \emph{Say it:} \emph{lājī}, once more
  \item \emph{Say it:} \emph{ménkǒu yǒu lājī} — There's rubbish by the door.
  \item \emph{Recall:} say the Mandarin for to go fishing, then the Mandarin for to mop the floor, then say \emph{lājī} again
\end{itemize}

\begin{tcolorbox}[breakable,colback=teal!4,colframe=teal!35!black,title={Before you move on}]
What does lājī mean? (\textquotedblleft{}Rubbish\textquotedblright{}.) Say it once more.
\end{tcolorbox}
\section[shài]{shài — to dry in the sun}

\label{lesson:ZH-C142-shai}



\begin{quote}
Before the new one: say the Mandarin for to mop the floor, then the Mandarin for rubbish.
\end{quote}

\subsection*{You'll want to know: shài}
\emph{shài} — \textquotedblleft{}to dry in the sun\textquotedblright{}.

\begin{quote}
\emph{jīntiān shì qíngtiān wǒmen shài yīfu} — It's a sunny day, so we're putting the washing out to dry.
\end{quote}

Its characters come later; for now it is a word for the ear.

\begin{cousinweb}
\emph{shài} is inherited from Middle Chinese \emph{sreaH} (Guangyun, in Baxter's transcription), glossed 'to sun, to expose to the sun'.
\end{cousinweb}

\begin{grammarlens}[title={What you've built}]
One more word for everyday actions.
\end{grammarlens}

\subsection*{Your turn}
\begin{itemize}
\raggedright
  \item \emph{Say it:} \emph{shài}
  \item \emph{Say it:} \emph{shài}, once more
  \item \emph{Say it:} \emph{jīntiān shì qíngtiān wǒmen shài yīfu} — It's a sunny day, so we're putting the washing out to dry.
  \item \emph{Recall:} say the Mandarin for to mop the floor, then the Mandarin for rubbish, then say \emph{shài} again
\end{itemize}

\begin{tcolorbox}[breakable,colback=teal!4,colframe=teal!35!black,title={Before you move on}]
What does shài mean? (\textquotedblleft{}To dry in the sun\textquotedblright{}.) Say it once more.
\end{tcolorbox}
\section[zhěnglǐ]{zhěnglǐ — to sort out; to put in order}

\label{lesson:ZH-C142-zhengli}



\begin{quote}
Before the new one: say the Mandarin for rubbish, then the Mandarin for to dry in the sun.
\end{quote}

\subsection*{You'll want to know: zhěnglǐ}
\emph{zhěnglǐ} — \textquotedblleft{}to sort out; to put in order\textquotedblright{}.

\begin{quote}
\emph{wǒ zhěnglǐ le wǒ de shū} — I've sorted out my books.
\end{quote}

Its characters come later; for now it is a word for the ear.

Two third tones: zhěnglǐ is said zhénglǐ, the first rising.

\begin{cousinweb}
\emph{zhěnglǐ} is built from \emph{zhěng} 'straight, even' and \emph{lǐ} 'to put in order', both inherited from Middle Chinese, \emph{tsyengX} and \emph{liX} (Guangyun, in Baxter's transcription); the dictionary quotes an older book for the second: to work jade.
\end{cousinweb}

\begin{grammarlens}[title={What you've built}]
One more word for everyday actions.
\end{grammarlens}

\subsection*{Your turn}
\begin{itemize}
\raggedright
  \item \emph{Say it:} \emph{zhěnglǐ}
  \item \emph{Say it:} \emph{zhěnglǐ}, once more
  \item \emph{Say it:} \emph{wǒ zhěnglǐ le wǒ de shū} — I've sorted out my books.
  \item \emph{Recall:} say the Mandarin for rubbish, then the Mandarin for to dry in the sun, then say \emph{zhěnglǐ} again
\end{itemize}

\begin{tcolorbox}[breakable,colback=teal!4,colframe=teal!35!black,title={Before you move on}]
What does zhěnglǐ mean? (\textquotedblleft{}To sort out; to put in order\textquotedblright{}.) Say it once more.
\end{tcolorbox}
\section[jiāohuā]{jiāohuā — to water the plants}

\label{lesson:ZH-C142-jiaohua}



\begin{quote}
Before the new one: say the Mandarin for to dry in the sun, then the Mandarin for to sort out.
\end{quote}

\subsection*{You'll want to know: jiāohuā}
\emph{jiāohuā} — \textquotedblleft{}to water the plants\textquotedblright{}.

\begin{quote}
\emph{jīntiān nǐ jiāohuā wǒ tuōdì} — Today you water the plants and I'll mop the floor.
\end{quote}

Its characters come later; for now it is a word for the ear.

\begin{cousinweb}
\emph{jiāohuā} is built from \emph{jiāo} 'to pour water on' and \emph{huā} 'flower', both inherited from Middle Chinese, \emph{kew} and \emph{xwae} (Guangyun, in Baxter's transcription).
\end{cousinweb}

\begin{grammarlens}[title={What you've built}]
One more word for everyday actions.
\end{grammarlens}

\subsection*{Your turn}
\begin{itemize}
\raggedright
  \item \emph{Say it:} \emph{jiāohuā}
  \item \emph{Say it:} \emph{jiāohuā}, once more
  \item \emph{Say it:} \emph{jīntiān nǐ jiāohuā wǒ tuōdì} — Today you water the plants and I'll mop the floor.
  \item \emph{Recall:} say the Mandarin for to dry in the sun, then the Mandarin for to sort out, then say \emph{jiāohuā} again
\end{itemize}

\begin{tcolorbox}[breakable,colback=teal!4,colframe=teal!35!black,title={Before you move on}]
What does jiāohuā mean? (\textquotedblleft{}To water the plants\textquotedblright{}.) Say it once more.
\end{tcolorbox}
